% !TEX root = ../main.tex

\frontchapter{Abstract}


Bulk Synchronous Parallel (BSP) workloads power large-scale scientific simulations and distributed
machine learning, running for long durations at massive scale where small inefficiencies translate
into major cost and time-to-discovery impacts.

Many of the most consequential inefficiencies arise from conditions unknowable before runtime,
including evolving data distributions, performance variability, and environment-dependent
pathologies.
This thesis argues that systematically addressing these bottlenecks requires feedback
loops that can reconstruct and act on summarized global views of distributed state while the job is
running.

We develop this claim through two studies.
On the data path, CARP is a streaming ingestion and
indexing system for checkpoint data that enables efficient range queries without expensive
post-processing.
CARP periodically reconstructs a view of global key distributions and uses it to
adapt range partitions as distributions evolve, matching the query performance of a full sort with
up to $4.9\times$ faster ingestion.
On the compute path, we study placement in adaptive mesh
refinement codes, where apparent load imbalance is confounded by fail-slow hardware, fabric
pathologies, and tuning artifacts.
Isolating true placement effects required progressively
eliminating each confounder---a process that consumed months even though individual mitigations
typically took days.
With reliable telemetry established, we design CPLX, a tunable placement
policy that achieves up to 21.6\% runtime improvement.
Across both studies, the decisive factor was
access to the right view at the right time.

Because these conditions recur unpredictably, production BSP systems need an always-on,
programmable observability loop: coarse views surface deviations, and the loop can be steered on
demand into narrower views to localize and act on them.
We call this steerable observability.
ORCA realizes this via a tree-based overlay for programmable feedback and control in BSP systems,
materializing runtime views through in-situ reduction and providing a timestep-consistent control
plane for coordinated interventions.
On a real AMR code at 4096 ranks, ORCA collects detailed
traces at 1--2\% overhead versus conventional HPC tracers at 56--81\% and eliminates 98.5--99.999\%
of telemetry at source for evaluated queries.
In an evaluation of its closed-loop capabilities
using an LLM agent, a performance anomaly that took months to diagnose manually is localized in 27
minutes.
